\chapter{Bodies That Change Shape}
\label{shape}

\section{The falling cat}

Drop a cat, upside down, from rest. In free fall there is no
external torque about the center of mass, so angular momentum is
conserved --- and it starts, and remains, at zero. Naive intuition
says that a body with identically zero angular momentum cannot
rotate. And yet the cat lands on its feet, having reoriented itself
by roughly $180^\circ$. The naive inference silently assumes a
\emph{rigid} body; it simply does not apply to a body that can
deform itself internally, and the cat's righting maneuver --- a
sequence of twists of its front and back halves relative to one
another --- is a real, controlled deformation, not a violation of
any conservation law.

\section{Shape space and the gauge structure of deformation}

The clean way to see why this works, due to Shapere and
Wilczek~\cite{shapere+wilczek} and to Montgomery's later "gauge
theory of the falling cat"~\cite{montgomery}, is to separate a
deformable body's configuration into its \emph{shape} $s$ (the
internal, orientation-independent degrees of freedom --- relative
joint angles, say) and an overall orientation $R\in SO(3)$ that fixes
that shape in space. Locally, the full configuration space is a
principal fiber bundle with structure group $SO(3)$ over shape space,
and the requirement that total angular momentum stay identically
zero throughout the motion is a linear relation between the shape
velocity $\dot s$ and the angular velocity $\omega$ of the body,
\begin{equation}
  \omega \;=\; -\,\mathcal A(s)\,\dot s ,
  \label{connection}
\end{equation}
where $\mathcal A(s)$ is built from the shape-dependent inertia
tensor and its couplings. This is precisely the data of a connection
(a gauge potential) on the bundle: given any path $s(t)$ in shape
space, \eqref{connection} lifts it uniquely to a path in the full
configuration space, exactly as a connection lifts a path in the
base of a fiber bundle to a horizontal path upstairs.

The consequence that matters here is that if the shape returns
exactly to where it started after a closed loop in shape space,
$s(T) = s(0)$, the orientation need \emph{not} return to where it
started. The net rotation accumulated around the loop is the
\emph{holonomy} of the connection~\eqref{connection}, and for a
small loop enclosing area $\Sigma$ it is governed, Stokes'-theorem
style, by the curvature $\mathcal F = d\mathcal A + \mathcal
A\wedge\mathcal A$ of that connection:
\begin{equation}
  \Delta R \;\approx\; \exp\!\left[\oint_{\partial\Sigma}\mathcal
  A\right] \;\approx\; \exp\!\left[\int_{\Sigma}\mathcal F\right].
  \label{holonomy}
\end{equation}
This is exactly the structure of a non-Abelian geometric (Berry,
Wilczek--Zee) phase from quantum mechanics, and of Aharonov--Bohm
holonomy in gauge theory --- reproduced in full by an entirely
classical, entirely mechanical system. It is also, viewed from the
vantage point of Section~\ref{cm:nh}, nothing but another
instance of a nonholonomic constraint: "total angular momentum
$\equiv 0$" is a Pfaffian relation between $\dot s$ and $\omega$ that
fails to be integrable to any function of shape and orientation
alone \emph{exactly when} the connection~\eqref{connection} has
nonzero curvature. The falling-cat effect and the rolling coin are,
formally, the same phenomenon in different clothing.

\section{Swimming without pushing}

The same structure governs self-propulsion in a highly viscous
fluid, where inertia is negligible and net displacement over one
cycle of body deformation is possible only if the cycle of shapes is
genuinely non-reciprocal. This is Purcell's scallop theorem: a
scallop that opens and closes through the \emph{same} sequence of
shapes, forwards then backwards, achieves zero net displacement,
because time-reversing the (inertia-free) equations of motion
exactly reverses the trajectory. Net swimming is again a
holonomy/geometric-phase effect in an appropriately defined shape
bundle~\cite{shapere+wilczek-swim}.

Perhaps the most striking illustration of the underlying idea is Jack
Wisdom's result that a quasi-rigid body executing a cyclic sequence
of internal shape changes in a curved \emph{spacetime} --- as
prescribed by general relativity --- can achieve a net translation
even in complete vacuum, with nothing to push against at
all~\cite{wisdom}. The effect scales with the local spacetime
curvature and with the extent of the body's internal motion; it is
minuscule for any realistic body and is not a practical propulsion
mechanism, but as a matter of principle it is "swimming" through
spacetime using exactly the same holonomy mechanism as the falling
cat, now realized by the genuine nonlinearity of general relativity
rather than by a gauge potential on a shape bundle constructed by
hand. Once again, an outcome that sounds like it must violate
momentum conservation turns out to be perfectly consistent with
conservation laws, provided one tracks correctly what "conservation"
does and does not forbid for a body that is allowed to change shape.
